\chapter{Myelodysplastic Syndromes}
\index{Myelodysplastic Syndromes}

\section*{Definition and Overview}
Myelodysplastic syndromes are clonal bone-marrow disorders in which blood-cell production is ineffective and abnormal. They cause anaemia, low platelets, low white cells, or combinations of these and may progress to acute myeloid leukaemia.

\section*{Assessment and Management}
Assessment is guided by the history, examination, and appropriate investigations. Management should be individualised by a qualified clinician and follow current Indian clinical guidance. Depending on the condition, care may include observation, medicines, procedures, surgery, rehabilitation, or specialist referral. Self-treatment should not delay evaluation of serious symptoms.

\section*{When to Seek Medical Care}
Fever, recurrent infection, unusual bleeding, or severe fatigue in a person with low blood counts requires prompt medical review.

\section*{India-specific Care}
In India, start with a qualified general physician or specialist at an appropriate clinic or hospital. Emergency symptoms should be taken to the nearest emergency department. Treatment availability varies by facility, so referral to a medical college or tertiary centre may be needed for complex disease.

\section*{Sources and Further Reading}
\begin{itemize}
    \item Indian Council of Medical Research (ICMR). Clinical and public-health guidance.
    \item Ministry of Health and Family Welfare, Government of India.
    \item World Health Organization. Relevant clinical and public-health information.
\end{itemize}

